\subsection{A tighter certificate forces full support but not product work}

Keep the path, exact-cell model, and every generation rule of
\cref{def:cert-prefix-poly}, but tighten only the common task to
\begin{equation}
 q:=\frac1n,
 \qquad
 \alpha_n=q^2,
 \qquad
 \epsppr^{\rm tight}:=\frac{q}{10}=\frac1{10n}.
 \label{eq:tight-certificate-path-parameters}
\end{equation}
Thus the task $\mathcal T_n^{\rm tight}$ accepts an explicit sparse list only
after an exact final verifier confirms
\begin{equation}
 \norm{D^{-1/2}(b-Q_n\widehat x)}_\infty
 \leq \alpha_n\epsppr^{\rm tight}=\frac{q^3}{10}.
 \label{eq:tight-sparse-path-residual-task}
\end{equation}
Unlisted coordinates are zero, there is one solve/output interaction, and the
verifier recomputes the residual from exposed path rows.  Write
$\mathsf{CertPrefixPoly}^{\rm tight}$ for the same executions as in
\cref{def:cert-prefix-poly} with this stricter task.  In particular, this
notation changes neither the allowed recurrence actions nor the exclusions.

Define the fixed spatial polynomial
\begin{equation}
 P(s):=\frac5{18}+\frac59s^2+\frac5{27}s^4+\frac2{81}s^6,
 \qquad 0\leq s\leq1.
 \label{eq:tight-profile-polynomial}
\end{equation}
Put $a:=(1+q^2)/2$, $\eta:=(1-q^2)/2$, and
\begin{equation}
 \gamma_n:=
 \frac{q}{aP(1-q)-\eta P(1-2q)},
 \qquad
 s_i:=\frac{n-i}{n}.
 \label{eq:tight-profile-normalizer}
\end{equation}
The candidate is specified in normalized coordinates by
\begin{equation}
 \widehat y_i:=\gamma_n qP(s_i),
 \qquad
 \widehat x_i:=\sqrt{d_i}\,\widehat y_i
 \quad(1\leq i\leq n).
 \label{eq:tight-profile-candidate}
\end{equation}
The degree six in \cref{eq:tight-profile-polynomial} is a spatial profile
degree.  As a polynomial in $Q_n$ applied to $b$, the final vector may have
degree $n-1$.

\begin{theorem}[A full-support sparse-basis recurrence beats the product
scale; \ProvedStatus]
\label{thm:tight-certificate-sparse-basis-counteralgorithm}
For every $n\geq8$:
\begin{enumerate}[label=(\roman*)]
\item every output accepted by $\mathcal T_n^{\rm tight}$ has all $n$
      coordinates nonzero, so its final support is $V(P_n)$ and
      $\nu_{\rm fin}=\nu_n=2(n-1)$;
\item the vector in \cref{eq:tight-profile-candidate} satisfies the strict
      version of \cref{eq:tight-sparse-path-residual-task} and is generated
      by an admissible $\mathsf{CertPrefixPoly}^{\rm tight}$ execution with
      \begin{equation}
       \mathcal C_{\rm sparse\text{-}basis}=\bigl(
       \Theta(n),n,1,0,\Theta(n),\Theta(n),0,
       \Theta(n),O(1),0,n\bigr);
       \label{eq:tight-sparse-basis-resource-vector}
      \end{equation}
\item on the identical task and chosen sequential adjacency interface, a
      fully charged append-only $LDL^\top$ response emits the exact solution
      and has
      \begin{equation}
       \mathcal C_{\rm tight\text{-}LDL}=\bigl(
       \Theta(n),n,1,0,\Theta(n),0,\Theta(n),
       \Theta(n),O(1),0,n\bigr).
       \label{eq:tight-path-response-resource-vector}
      \end{equation}
\end{enumerate}
Both executions are deterministic exact-algebraic-cell algorithms, with no
preprocessing, randomness, or failure probability.  Thus their charged
nonmemory work is
$\Theta(n)=\Theta(\nu_{\rm fin})$, whereas
\begin{equation}
 \frac{\nu_{\rm fin}}{\sqrt{\alpha_n}}=\Theta(n^2).
 \label{eq:tight-path-product-refutation}
\end{equation}
The stricter accuracy therefore does not restore a product lower bound for
the broad supported-prefix polynomial semantics.
\end{theorem}

\begin{proof}
We first prove the full-support assertion.  Let
$y^0:=D^{-1/2}x^0$ and
\begin{equation*}
 \chi:=\frac{1-q}{1+q}.
\end{equation*}
Solving the second-order path recurrence with its two ambient endpoint rows
gives
\begin{equation}
 y_i^0
 =q\frac{\chi^{i-1}+\chi^{2n-i-1}}
          {1-\chi^{2n-2}}
 \quad(1\leq i\leq n).
 \label{eq:tight-path-exact-profile}
\end{equation}
The numerator decreases through $i=n$.  The companion exact-CG calibration
proves $\chi^{-n}\leq8$ for $n\geq3$, and hence
\begin{equation}
 y_i^0\geq y_n^0
 =\frac{2q\chi^{n-1}}{1-\chi^{2n-2}}
 >2q\chi^{n-1}\geq\frac q4.
 \label{eq:tight-path-minimum-solution}
\end{equation}
For any residual $r=b-Q_n\widehat x$, the matrix
$D^{-1/2}Q_n^{-1}D^{1/2}$ is nonnegative and has every row sum equal to
$1/\alpha_n$: this follows from the nonsingular $M$-matrix property and
$Q_n\sqrt d=\alpha_n\sqrt d$.  Therefore the certificate implies
\begin{equation*}
 \norm{D^{-1/2}(x^0-\widehat x)}_\infty
 \leq\alpha_n^{-1}\norm{D^{-1/2}r}_\infty
 \leq\frac q{10}.
\end{equation*}
If any coordinate were unlisted, its normalized error would be
$y_i^0>q/4$, a contradiction.  This proves (i), including the value of
$\nu_{\rm fin}$.

We next check the candidate residual exactly.  Since $P$ is increasing and
\begin{equation*}
 P'(3/4)=\frac{907}{768},
\end{equation*}
the denominator in \cref{eq:tight-profile-normalizer} obeys, for $q\leq1/8$,
\begin{align*}
 aP(1-q)-\eta P(1-2q)
 &=\frac{P(1-q)-P(1-2q)}2
   +\frac{q^2}{2}\bigl(P(1-q)+P(1-2q)\bigr)\\
 &\geq\frac q2P'(3/4)=\frac{907}{1536}q.
\end{align*}
Thus $\gamma_n$ is positive and
\begin{equation}
 0<\gamma_n\leq\frac{1536}{907}.
 \label{eq:tight-profile-normalizer-bound}
\end{equation}
Let $h:=D^{-1/2}(b-Q_n\widehat x)$.  The choice of $\gamma_n$ gives
$h_1=0$.  For $2\leq i\leq n-1$, the exact centered-difference identity for
the degree-six polynomial gives
\begin{equation}
 h_i=\gamma_n q^3R_q(s_i),
 \label{eq:tight-profile-interior-residual}
\end{equation}
where
\begin{equation}
 R_q(s):=-P(s)+\frac{1-q^2}{4}P''(s)
       +\frac{(1-q^2)q^2}{48}P^{(4)}(s)
       +\frac{(1-q^2)q^4}{1440}P^{(6)}(s).
 \label{eq:tight-profile-residual-polynomial}
\end{equation}
The coefficients in \cref{eq:tight-profile-polynomial} were chosen so that
\begin{equation*}
 \frac14P''(s)-P(s)=-\frac2{81}s^6.
\end{equation*}
Moreover,
$P''(1)=110/27$, $P^{(4)}(1)=40/3$, and $P^{(6)}=160/9$;
these derivatives are nonnegative and increasing where applicable on
$[0,1]$.  Rewriting \cref{eq:tight-profile-residual-polynomial} around the
preceding cancellation gives
\begin{equation*}
 R_q(s)=\left(\frac14P''(s)-P(s)\right)
       -\frac{q^2}{4}P''(s)
       +\frac{(1-q^2)q^2}{48}P^{(4)}(s)
       +\frac{(1-q^2)q^4}{1440}P^{(6)}(s).
\end{equation*}
Therefore the triangle inequality, $1-q^2\leq1$, and the displayed
derivative maxima give, for $q\leq1/8$,
\begin{align}
 |R_q(s)|
 &\leq\frac2{81}
   +\left(\frac{55}{54}+\frac5{18}\right)q^2
   +\frac1{81}q^4
 \nonumber\\
 &=\frac2{81}+\frac{35}{27}q^2+\frac1{81}q^4
 \leq\frac{1657}{36864},
 \nonumber\\
 \frac{|h_i|}{q^3}
 &\leq\frac{1536}{907}\frac{1657}{36864}
 =\frac{1657}{21768}<\frac1{10}.
 \label{eq:tight-profile-interior-bound}
\end{align}
At the far endpoint, direct substitution gives
\begin{equation}
 h_n=-\gamma_n q^5
 \left(\frac5{27}+\frac{13}{162}q^2+\frac1{81}q^4\right),
 \label{eq:tight-profile-endpoint-residual}
\end{equation}
and therefore
\begin{equation}
 \frac{|h_n|}{q^3}
 \leq\frac{2291}{464384}<\frac1{10}.
 \label{eq:tight-profile-endpoint-bound}
\end{equation}
Together with $h_1=0$, this proves the strict certificate.

It remains to realize the candidate without a hidden response.  Put
$c_j:=-(Q_n)_{j,j+1}>0$.  Start with
$u_1:=b/\alpha_n=e_1$.  After exposing the next path row, use one charged
sparse row action and the exact scalar combination
\begin{equation}
 u_{j+1}:=
 \frac{a u_j-c_{j-1}u_{j-1}-Q_nu_j}{c_j}=e_{j+1}
 \quad(1\leq j<n),
 \label{eq:tight-coordinate-basis-recurrence}
\end{equation}
where the $c_0u_0$ term is omitted.  Inductively, $u_j=e_j$ is a polynomial
of degree $j-1$ in $Q_n$ applied to $b$, stored on a singleton inside its
allowed reached prefix.  The cancellations in
\cref{eq:tight-coordinate-basis-recurrence} are exact scalar linear
combinations, not interior masks.  Every row action touches only a constant
number of path entries.  Once the far endpoint reveals $n$, evaluate the
fixed formula
$\sqrt{d_i}\gamma_nqP(s_i)$ and combine it with $u_i$ for each coordinate.
Each exposed $\sqrt{d_i}$ is merely the scalar coefficient of one singleton
direction; no non-scalar diagonal map is applied.
This delayed synthesis uses $\Theta(n)$ charged recurrence arithmetic.

The chosen implementation exposes and caches the path in $n$ sequential
adjacency batches, stores its rows and the singleton directions, and uses a
constant-size rolling workspace.  A first streamed profile pass recomputes
the exact residual from cached rows; after it succeeds, a second streamed pass
emits the $n$ cells.  Coefficient evaluation and scalar combinations belong
to $C_{\rm rec}$, while loop, stopping, and verifier arithmetic belong to
$C_{\rm ctl}$.  No complete current active vector is maintained, so
$C_{\rm mat}=0$.  This proves every coordinate of
\cref{eq:tight-sparse-basis-resource-vector}; $R_{\rm adj}=n$ is the cost of
this nonprefetching realization, not an adaptivity lower bound.

For the response arm, use the same sequential exposure to append every
tridiagonal pivot, multiplier, and forward message.  Two linear reverse
passes suffice to recompute and verify the exact solution and then emit its
$n$ nonzero cells, using constant rolling workspace and no intermediate
complete vector.  Every factor/message construction and reverse application
is charged to $C_{\rm resp}$, and the residual scan is charged to
$C_{\rm ctl}$.  This proves
\cref{eq:tight-path-response-resource-vector}.  The work comparison and
\cref{eq:tight-path-product-refutation} now follow from
$\nu_{\rm fin}=2(n-1)$ and $\sqrt{\alpha_n}=1/n$.
\end{proof}

\paragraph{Classification of sparse propagation and delayed synthesis.}
The coordinate-basis construction is sparse Lanczos propagation specialized
to an unreduced path: exposed-row actions and scalar cancellations generate
its directions, so their arithmetic is $C_{\rm rec}$.  Finite propagation
still forces $n-1$ such actions before the far coordinate is available, but
each generated direction is a singleton rather than a full prefix.  The fixed
profile in
\cref{eq:tight-profile-polynomial} is evaluated only after $n$ is exposed;
its numerical combinations remain $C_{\rm rec}$ and its scheduling is
$C_{\rm ctl}$.  Its universal rational coefficients are not graph-dependent
preprocessing, and no inverse coefficients are precomputed.  Any
graph-dependent coefficient table prepared before the online phase would be
reported in $C_{\rm pre}$ and its response representation.  By contrast, if an
implementation builds a projected tridiagonal system and computes its
coefficients with Thomas pivots, a factor, a Green formula, or another
graph-dependent inverse map, that construction and application are
$C_{\rm resp}$, with its state and scratch charged to the two memory
coordinates.  Calling such a scalar solve ``coefficient synthesis'' would
not make it free recurrence work.  Conversely, the class does not require
full-prefix scans when exact cancellations leave singleton directions; an
implementation that chooses to maintain a complete current prefix would pay
the corresponding $C_{\rm mat}$ charge.

\paragraph{Exact scope of the tighter refutation.}
The tighter task removes the constant-prefix escape: every certified output
has full support and therefore pays the optimal $\Theta(n)$ exposure and
emission scale.  It still leaves a linear-work supported recurrence because
the broad semantics permit sparse coordinate directions and delayed scalar
combinations.  The result is exact-cell and path-specific.  It neither
shortens ordinary CG, supplies a finite-precision implementation, refutes the
exact full-vector $\mathsf{DiagSpecPoly}(r)$ theorem, nor rules out a product
lower bound for a narrower independently motivated materialization or
trajectory class or for a different graph family.
